\booksection{第一章付録：HarmonyOSオペレーティングシステム特集}

\chapterintro{HarmonyOSはHuaweiが自主開発したオペレーティングシステムで、2019年の発表から現在までにいくつかのメジャーバージョンを経てきた。それは単なる「もう一つのAndroid代替品」ではなく、独特な設計理念を持つ新種だ。この台頭しつつあるオペレーティングシステムを一緒に深く理解しよう。}

\booksubsection{一、HarmonyOSとは？なぜ他と違うのか？}

HarmonyOSの核心理念は「万物互聯」。従来のオペレーティングシステムが主に単一デバイスにサービスを提供するのとは異なり、HarmonyOSは設計当初からマルチデバイス連携を考慮している。

\textbf{分散アーキテクチャ}はHarmonyOSの魂だ。簡単に言うと、複数のデバイスを仮想的に1つの「スーパーターミナル」にできる。あなたのスマートフォン、タブレット、パソコン、時計、テレビはもはや孤立した個体ではなく、1つの全体の「器官」となり、互いの能力を呼び出し合える。

例えば：スマートフォンで動画を見ていて大画面で見たくなった時、従来の方法はスクリーンミラーリングだが、HarmonyOSでは動画を直接テレビに「流転」して続きを見られ、スマートフォンはリモコンになる。これは単なるスクリーンミラーリングではなく、アプリ自体がデバイスをまたいで実行されているのだ。

\booksubsection{二、HarmonyOS 4.3：AOSPベースの成熟バージョン}

HarmonyOS 4.3は現在Huaweiスマートフォンに広くデプロイされているバージョンだ。説明が必要なのは、HarmonyOS 4.3はAndroidオープンソースプロジェクト（AOSP）のコードフレームワークと互換性があるため、Androidアプリを完全に互換実行できる——Huaweiアプリマーケットでダウンロードする大部分のアプリは本質的に依然としてAndroidアプリで、HarmonyOSの互換レイヤーを通じて実行されている。しかしHarmonyOS 4.3は単なるAndroidの皮ではなく、AOSPをベースにHuawei自社開発の分散能力、方舟エンジンなどの技術を大量に追加している。

\subsubsection*{コア特性：}

\begin{itemize}
  \item \textbf{Androidアプリ互換}：AOSPベースで開発されているため、HarmonyOS 4.3はほぼすべてのAndroidアプリを正常に実行でき、ユーザーはほとんど差を感じない。
  \item \textbf{方舟エンジン}：Huawei自社開発のコンパイラとランタイムで、アプリの実行をより滑らかにより省電力にする。早期バージョンより滑らかさが明確に向上。
  \item \textbf{スーパーターミナル}：コントロールセンターで周囲の他のHarmonyOSデバイスを直感的に確認でき、ドラッグで接続。スマートフォンがパソコンのキーボードマウスを呼び出したり、タブレットがスマートフォンのカメラを呼び出したりするのが便利。
  \item \textbf{サービスカード}：iOSのウィジェットに似ているがより柔軟。アプリアイコンを上スワイプしてカードを展開し、アプリを開かずに主要情報を表示。
  \item \textbf{または}：公式アプリマーケットからのアプリインストールのみを許可し、マルウェアリスクを低減。
  \item \textbf{プライバシーセンター}：アプリ権限を一元管理し、アプリの行動記録を確認。
\end{itemize}

\subsubsection*{ユーザーインターフェースの特徴：}

\begin{itemize}
  \item \textbf{マガジンロック画面}：ロック画面に美しい画像と文章を表示し、スワイプするとデスクトップに入る。
  \item \textbf{万能カード}：デスクトップアイコンをカードに展開し、より多くの情報を表示。
  \item \textbf{スーパー中継ステーション}：一時保存領域で、文字、画像、ファイルを一時保存し、任意のアプリにドラッグして使用。アプリ間共有が極めてスムーズに。
  \item \textbf{スマート検索}：グローバル検索でアプリ、ファイル、設定、ウェブページを検索。
\end{itemize}

\booksubsection{三、HarmonyOS NEXT（HarmonyOS 6）：Huawei自主オペレーティングシステムの真の起点}

HarmonyOS NEXT（HarmonyOS NEXT 5.0）はHarmonyOS 4.3とは根本的に異なる——Huawei自社開発のHarmonyOSマイクロカーネルを採用し、LinuxカーネルとAOSPコードに依存せず、\textbf{Androidアプリを直接サポートしない}。これはHarmonyOSが「Android互換の過渡形態」から「完全に独立したオペレーティングシステム」へ正式に移行する重要な一歩だ。

\subsubsection*{重大な変革：}

\begin{itemize}
  \item \textbf{完全自主カーネル}：HarmonyOS NEXTはHuawei自社開発のカーネルを使用し、AOSPコードを一切含まない。システムはより簡素で、より安全で、より効率的。
  \item \textbf{Androidアプリ非サポート}：これが最大の変化。HarmonyOS NEXTはネイティブHarmonyOSアプリしか実行できず、APKファイルのインストールと実行は不可。アップグレード前に常用アプリにHarmonyOSネイティブ版があるか確認すること。
  \item \textbf{ネイティブアプリエコシステム}：Huaweiは開発者がアプリをHarmonyOSネイティブプラットフォームに移行することを全力で推進。現在WeChat、Alipay、Taobaoなどの主流アプリはHarmonyOSネイティブ版をリリース済みだが、一部のニッチアプリはまだ適応されていない可能性がある。
  \item \textbf{星盾セキュリティアーキテクチャ}：全く新しいセキュリティ体系で、ハードウェアからソフトウェアまでの全リンク保護。より細粒度の権限制御をサポート、例えば「あいまい測位」——アプリはあなたの大体の位置しか分からず、正確な位置は分からない。
  \item \textbf{AI深融合}：システムレベルのAI能力で、単なる特定のアプリにAIがあるのではない。例えばAI字幕は動画をリアルタイム翻訳し、AI切り抜きは主体をスマート認識。
  \item \textbf{デバイス間シームレス移行}：さらに進んだ分散能力。アプリ状態はデバイス間でシームレスに切り替えられ、スマートフォンで編集したドキュメントがパソコンの前に来ると自動的にパソコンで開く。
  \item \textbf{元サービス}：インストール不要で使用可能な軽量サービス。例えばQRコードをスキャンして直接特定機能を開き、アプリを先にダウンロードする必要がない。
\end{itemize}

\begin{tcbwarn}{HarmonyOS NEXTアップグレード前の必読事項}
HarmonyOS NEXTにアップグレードすると、スマートフォンはAndroidアプリを実行できなくなる。アップグレード前に必ず確認すること：(1)常用アプリにHarmonyOSネイティブ版があるか；(2)仕事や生活にAndroidアプリに依存する必須シーンがないか。アップグレード後はHarmonyOS 4.3に戻すことはできない。

\end{tcbwarn}

\subsubsection*{操作ロジックの変化：}

\begin{itemize}
  \item \textbf{ジェスチャーナビゲーション}：左右のエッジスワイプで戻る、下部上スワイプでデスクトップへ、上スワイプ長押しでマルチタスクへ。
  \item \textbf{スマートマルチウィンドウ}：分割画面、フローティングウィンドウ、平行視界など複数のマルチタスク形態をサポート。
  \item \textbf{Changlian通話}：Huaweiのビデオ通話サービスで、画面共有、ビューティー、背景ボケをサポート。
\end{itemize}

\booksubsection{四、HarmonyOSエコシステム：デバイスとアプリ}

HarmonyOSは単なるスマートフォンシステムではなく、多種デバイスをカバーするエコシステムだ。

\subsubsection*{サポートするデバイスタイプ：}

\begin{itemize}
  \item \textbf{スマートフォンとタブレット}：コアデバイスで機能が最も完全。
  \item \textbf{スマートウォッチとバンド}：単独実行またはスマートフォンと連携。
  \item \textbf{スマートスクリーン（テレビ）}：家庭エンターテインメントの中心。
  \item \textbf{スマートスピーカー}：音声インタラクションの入り口。
  \item \textbf{自動車}：HarmonyOSコックピット、スマート運転席体験。
  \item \textbf{スマートホーム}：冷蔵庫、エアコン、洗濯機など。
\end{itemize}

\subsubsection*{アプリエコシステムの現状：}

HarmonyOSアプリマーケットには数百万のアプリがあり、主要なアプリはほぼ揃っています。Huaweiも開発者によるネイティブHarmonyOSアプリの開発を積極的に推進しています。2026年中盤時点で、WeChat、Alipay、Taobao、Douyin、Meituan、JD.comなどの主要アプリはすべてHarmonyOSネイティブ版をリリースしており、日常使用はほぼ問題ありません。一部のニッチアプリや業界向けアプリはまだ適配中の可能性があるため、アップグレード前に常用アプリが対応しているか確認することをお勧めします。

\booksubsection{五、HarmonyOS実用テクニック}

\begin{itemize}
  \item \textbf{スーパーターミナル}：画面右上から下スワイプしてコントロールセンターを開き、「スーパーターミナル」をタップすると周囲の他のデバイスが表示され、ドラッグで接続。
  \item \textbf{スーパー中継ステーション}：文字/画像/ファイルを長押し、「中継ステーションにコピー」を選択し、他のアプリで長押しペーストまたはドラッグで取り出す。
  \item \textbf{サービスカード}：アプリアイコンを上スワイプしてカードを展開、カードを長押しでサイズ調整やデスクトップに追加。
  \item \textbf{スマートマルチウィンドウ}：画面の左右エッジから内側にスワイプして長押し、サイドバーを開き、アプリをタップしてフローティングウィンドウで開く。
  \item \textbf{プライバシー保護}：設定 > プライバシーで、アプリ権限使用記録の確認、あいまい測位の設定などが可能。
\end{itemize}

\begin{tcbwarn}{注意事項}
HarmonyOSシステムは絶えず更新反復されており、バージョンにより多少の差異がある可能性がある。本記事の記述はHarmonyOS 4.3と6の主要特性に基づいており、具体的な機能はあなたのスマートフォンの実際のバージョンによる。システム更新を維持し、最新機能とセキュリティパッチを取得することを推奨。

\end{tcbwarn}

\begin{tcbpractice}{今すぐ試そう：HarmonyOSの特色機能を体験}
\begin{enumerate}
  \item コントロールセンターを開き、「スーパーターミナル」を見つけ、周囲に他のHarmonyOSデバイスがあるか確認。
  \item 「スーパー中継ステーション」を試す：文字を中継ステーションにコピーし、メモ帳でドラッグして取り出す。
  \item 常用アプリのアイコンをいくつか上スワイプし、サービスカードが何を表示するか見てみる。
  \item 設定 > プライバシーで、最近どのアプリがセンシティブな権限を使用したか確認。
\end{enumerate}

\end{tcbpractice}
